\documentclass[12pt, a4paper]{article}
\usepackage[utf8]{inputenc}
\usepackage[T1]{fontenc}
\usepackage[ngerman]{babel}
\usepackage{amsmath, amssymb, amsthm}
\usepackage{geometry}
\usepackage{graphicx}
\usepackage{xcolor}
\usepackage{tcolorbox}
\usepackage{enumitem}
\usepackage{listings}
\usepackage{fancyhdr}
\usepackage{titlesec}
\usepackage{mdframed}
\usepackage{tikz}
\usepackage{pgfplots}
\pgfplotsset{compat=1.18}

\geometry{a4paper, margin=2.5cm}

\definecolor{theoryblue}{RGB}{30, 90, 160}
\definecolor{exerciseorange}{RGB}{200, 80, 0}
\definecolor{solutiongreen}{RGB}{0, 120, 60}
\definecolor{hintgray}{RGB}{100, 100, 100}
\definecolor{codebg}{RGB}{245, 245, 245}

% Theory box
\tcbuselibrary{skins, breakable}
\newtcolorbox{theorybox}[1]{
  colback=theoryblue!5,
  colframe=theoryblue,
  fonttitle=\bfseries\color{white},
  title=#1,
  breakable
}

% Exercise box
\newtcolorbox{exercisebox}[1]{
  colback=exerciseorange!5,
  colframe=exerciseorange,
  fonttitle=\bfseries\color{white},
  title=Aufgabe #1,
  breakable
}

% Solution space box
\newtcolorbox{solutionspace}[1][6cm]{
  colback=white,
  colframe=black!20,
  fonttitle=\itshape\color{hintgray},
  title=Rechenraum,
  height=#1,
  breakable
}

% Hint box
\newtcolorbox{hintbox}{
  colback=solutiongreen!5,
  colframe=solutiongreen,
  fonttitle=\bfseries\color{white},
  title=Hinweis,
  breakable
}

% Julia code style
\lstdefinelanguage{Julia}{
  keywords={function, end, return, if, else, elseif, for, while, in, begin, let, do, try, catch, finally, using, import, export, mutable, struct, abstract, type},
  keywordstyle=\color{theoryblue}\bfseries,
  comment=[l]{\#},
  commentstyle=\color{hintgray}\itshape,
  stringstyle=\color{solutiongreen},
  morestring=[b]",
  sensitive=true
}
\lstset{
  language=Julia,
  backgroundcolor=\color{codebg},
  basicstyle=\ttfamily\small,
  breaklines=true,
  frame=single,
  rulecolor=\color{black!30},
  numbers=left,
  numberstyle=\tiny\color{hintgray}
}

\pagestyle{fancy}
\fancyhf{}
\rhead{\textcolor{hintgray}{\small Rocket Science Workbook}}
\lhead{\textcolor{hintgray}{\small \leftmark}}
\cfoot{\thepage}

\titleformat{\section}{\large\bfseries\color{theoryblue}}{}{0em}{}[\titlerule]
\titleformat{\subsection}{\normalsize\bfseries\color{exerciseorange}}{}{0em}{}


\begin{document}

\begin{center}
  {\LARGE \textbf{Chapter 06: Multi-Stage Rockets}}\\[0.5em]
  {\large \textcolor{hintgray}{Staging Equations, Mass Optimization, and Discontinuous ODEs}}\\[0.3em]
  {\small \textcolor{hintgray}{Rocket Science Workbook — simulation-2D-jl}}
\end{center}
\vspace{1em}\hrule\vspace{1.5em}

% ──────────────────────────────────────────────────────────────
\section{Why Staging is Necessary}

\begin{theorybox}{The Mass Fraction Problem}
For a single-stage rocket to reach LEO, we need $\Delta v \approx 9400\;\text{m/s}$.
With $v_e = 3200\;\text{m/s}$ (RP-1/LOX):
\[
  \frac{m_0}{m_f} = e^{9400/3200} = e^{2.9375} \approx 18.9
\]
So 94.7\% of the launch mass must be propellant. The structural fraction (tanks, engines, plumbing) for real rockets is at least 5–10\% of total mass.\\[0.4em]
Structural fraction $\varepsilon = 5\%$ of propellant mass $\Rightarrow$ only $\approx 0.3\%$ is payload.\\[0.4em]
\textbf{Staging solves this}: discard empty tanks to stop carrying dead mass.
\end{theorybox}

% ──────────────────────────────────────────────────────────────
\section{Multi-Stage Rocket Equation}

\begin{theorybox}{$\Delta v$ for $n$ Stages}
Each stage fires independently and is then jettisoned. The total $\Delta v$ is additive:
\[
  \Delta v_\text{total} = \sum_{i=1}^n v_{e,i}\,\ln\frac{m_{0,i}}{m_{f,i}}
\]
where $m_{0,i}$ includes all stages above (plus payload) at the start of stage $i$,
and $m_{f,i}$ excludes the dry mass of stage $i$ just before separation.
\end{theorybox}

\subsection*{Structural Coefficient}

\begin{theorybox}{Definition}
\[
  \varepsilon_i = \frac{m_{\text{dry},i}}{m_{\text{wet},i}} = \frac{\text{empty stage mass}}{\text{full stage mass including fuel}}
\]
Typical values: $\varepsilon \approx 0.05$–$0.12$.

The mass ratio for stage $i$ alone:
\[
  \mu_i = \frac{m_{\text{wet},i}}{m_{\text{dry},i}} = \frac{1}{\varepsilon_i}
\]
Stage $\Delta v$ alone: $\Delta v_i = v_e\ln\mu_i = v_e\ln(1/\varepsilon_i)$
\end{theorybox}

% ──────────────────────────────────────────────────────────────
\section{Optimal Staging}

\begin{theorybox}{Equal-Stage Optimality}
For $n$ stages with identical $v_e$ and $\varepsilon$, the maximum total $\Delta v$ is achieved when all stages are geometrically similar (same mass ratio). The total $\Delta v$:
\[
  \Delta v_\text{total} = n\cdot v_e\,\ln\!\left(\frac{1}{\varepsilon}\right)
\]
This scales \textit{linearly} with $n$ — but diminishing returns apply because of the overhead of additional stages (complexity, separation mechanisms, interstage mass).
\end{theorybox}

% ──────────────────────────────────────────────────────────────
\section{Discontinuities in the ODE}

\begin{theorybox}{Stage Separation as an Event}
At separation time $t^*$, the mass drops instantaneously:
\[
  m(t^{*-}) = m(t^{*+}) + m_{\text{dry, stage}}
\]
This is a \textbf{discontinuity} in the state — the ODE changes character at $t^*$.

In numerical integration: after every RK4 step, check if the active stage's fuel $\leq 0$.
If yes:
\begin{enumerate}
  \item Subtract $m_\text{dry}$ from the current mass
  \item Switch to next stage parameters ($T$, $\dot{m}$, $A$)
  \item Continue integration with new parameters
\end{enumerate}
\textbf{Important:} Do not ``straddle'' the discontinuity with a RK4 step — it corrupts the solution.
Ideally, stop exactly at $m_\text{fuel} = 0$ using event detection (bisection on $t^*$).
\end{theorybox}

% ──────────────────────────────────────────────────────────────
\section{Exercises}

\begin{exercisebox}{1 — Full Tsiolkovsky Derivation for 2 Stages}
A 2-stage rocket with payload $m_L = 13\,150\;\text{kg}$.\\
Stage 2: $m_{\text{wet},2} = 92\,670\;\text{kg}$, $\varepsilon_2 = 0.043$, $v_{e,2} = 3\,430\;\text{m/s}$.\\
Stage 1: $m_{\text{wet},1} = 418\,930\;\text{kg}$, $\varepsilon_1 = 0.051$, $v_{e,1} = 3\,050\;\text{m/s}$.

\textbf{a)} Compute $m_{\text{dry},2} = \varepsilon_2 \cdot m_{\text{wet},2}$:
\[
  m_{\text{dry},2} = 0.043 \times 92\,670 = \underline{\hspace{3cm}}\;\text{kg}
\]

\textbf{b)} At the start of stage 2, the total mass is $m_{\text{wet},2} + m_L$. At burnout, it is $m_{\text{dry},2} + m_L$.
\[
  \Delta v_2 = 3430\cdot\ln\!\left(\frac{92\,670 + 13\,150}{m_{\text{dry},2} + 13\,150}\right) = \underline{\hspace{5cm}}\;\text{m/s}
\]

\textbf{c)} At the start of stage 1, total mass is $m_{\text{wet},1} + m_{\text{wet},2} + m_L$.
At burnout (just before separation), it is $m_{\text{dry},1} + m_{\text{wet},2} + m_L$.
\[
  \Delta v_1 = 3050\cdot\ln\!\left(\frac{418\,930 + 92\,670 + 13\,150}{m_{\text{dry},1} + 92\,670 + 13\,150}\right) = \underline{\hspace{3cm}}\;\text{m/s}
\]

\textbf{d)} Total ideal $\Delta v$: \underline{\hspace{3cm}} m/s. Compare to the target of $9400\;\text{m/s}$.
\end{exercisebox}

\begin{solutionspace}[10cm]
\end{solutionspace}

\begin{exercisebox}{2 — Separation Timing}
Stage 1: $m_{\text{wet},1} = 418\,930\;\text{kg}$, $m_{\text{dry},1} = 21\,200\;\text{kg}$, $\dot{m}_1 = 2530\;\text{kg/s}$.

\textbf{a)} Available propellant in stage 1:
\[
  m_{\text{prop},1} = m_{\text{wet},1} - m_{\text{dry},1} = \underline{\hspace{3cm}}\;\text{kg}
\]

\textbf{b)} Burnout time (time from ignition until stage 1 propellant is exhausted):
\[
  t_\text{MECO1} = \frac{m_{\text{prop},1}}{\dot{m}_1} = \underline{\hspace{3cm}}\;\text{s}
\]

\textbf{c)} Just before separation, the total rocket mass is:
\[
  m(t_\text{MECO1}^-) = m_{\text{dry},1} + m_{\text{wet},2} + m_L = \underline{\hspace{3cm}}\;\text{kg}
\]
Just after (stage 1 dry mass jettisoned):
\[
  m(t_\text{MECO1}^+) = m_{\text{wet},2} + m_L = \underline{\hspace{3cm}}\;\text{kg}
\]

\textbf{d)} What fraction of launch mass has been discarded by this point?
\[
  \frac{m_0 - m(t_\text{MECO1}^+)}{m_0} = \underline{\hspace{3cm}}
\]
\end{exercisebox}

\begin{solutionspace}[7cm]
\end{solutionspace}

\begin{exercisebox}{3 — Staging Benefit: 1 vs 2 vs 3 Stages}
Same total $\Delta v_\text{target} = 9400\;\text{m/s}$, $v_e = 3200\;\text{m/s}$, $m_L = 13\,150\;\text{kg}$, $\varepsilon = 0.06$.

For $n$ equal stages, the payload fraction is:
\[
  \lambda = \frac{m_L}{m_0} = \left[\varepsilon + (1-\varepsilon)\,e^{-\Delta v / (n\,v_e)}\right]^n
\]

\textbf{a)} For $n=1$:
\[
  \lambda = \varepsilon + (1-\varepsilon)\,e^{-9400/3200} = \underline{\hspace{4cm}}
\]
Launch mass: $m_0 = m_L / \lambda = \underline{\hspace{3cm}}\;\text{kg}$

\textbf{b)} For $n=2$ (each stage has $\Delta v/2 = 4700\;\text{m/s}$):
\[
  \lambda = \left[\varepsilon + (1-\varepsilon)\,e^{-4700/3200}\right]^2 = \underline{\hspace{4cm}}
\]
Launch mass: $m_0 = \underline{\hspace{3cm}}\;\text{kg}$

\textbf{c)} For $n=3$ (each $\Delta v/3 \approx 3133\;\text{m/s}$):
\[
  \lambda = \left[\varepsilon + (1-\varepsilon)\,e^{-3133/3200}\right]^3 = \underline{\hspace{4cm}}
\]

\textbf{d)} Tabulate and discuss the diminishing returns of adding stages:
\begin{center}
\begin{tabular}{|c|c|c|c|}
\hline
Stages $n$ & Payload fraction $\lambda$ & Launch mass $m_0$ (kg) & Reduction vs $n-1$ \\
\hline
1 & & & — \\
\hline
2 & & & \\
\hline
3 & & & \\
\hline
\end{tabular}
\end{center}
\end{exercisebox}

\begin{solutionspace}[10cm]
\end{solutionspace}

\begin{exercisebox}{4 — Velocity at Stage Separation: Energy Conservation}
Just before separation, the rocket has velocity $|\mathbf{v}| = v_\text{sep}$ and altitude $h_\text{sep}$.
The stage 1 casing (mass $m_{\text{dry},1}$) becomes a ballistic projectile.

\textbf{a)} Given $v_\text{sep} = 2200\;\text{m/s}$, $h_\text{sep} = 70\;\text{km}$, $\theta_\text{sep} = 30°$ (from horizontal):
Find the horizontal and vertical velocity components at separation.

\textbf{b)} Using energy conservation (ignore drag), find the maximum altitude the stage casing reaches:
\[
  h_\text{max} = h_\text{sep} + \frac{v_y^2}{2g} = \underline{\hspace{4cm}}\;\text{km}
\]

\textbf{c)} Estimate the downrange distance where the stage casing lands (ballistic trajectory from $h_\text{sep}$).
\textit{Hint: Time of flight from $h_\text{sep}$ back to sea level, then multiply by $v_x$.}
\end{exercisebox}

\begin{solutionspace}[9cm]
\end{solutionspace}

% ──────────────────────────────────────────────────────────────
\section{Resources}

\begin{hintbox}
\textbf{Videos:}
\begin{itemize}
  \item \textbf{Everyday Astronaut} — ``Why multi-staging?'' (very accessible):\\
        \texttt{youtube.com/c/EverydayAstronaut}
  \item \textbf{Scott Manley} — ``How Falcon 9 stages work'' (engineering detail):\\
        Search: ``Scott Manley Falcon 9 staging''
  \item \textbf{MIT OpenCourseWare} — 16.06 ``Principles of Automatic Control'' Lecture 2\\
        (covers event detection and hybrid systems, relevant for staging logic)
\end{itemize}
\textbf{Reference:}
\begin{itemize}
  \item Curtis, \textit{Orbital Mechanics for Engineering Students}, Ch.\ 1 (rocket staging)
  \item \texttt{braeunig.us/space/rocket.htm} — Rocket staging calculations
\end{itemize}
\end{hintbox}

\section{Connection to the Julia Template}

\begin{theorybox}{What you implement in \texttt{template.jl}}
\begin{enumerate}
  \item \texttt{Stage} struct with all physical parameters
  \item \texttt{simulate\_multistage} with post-step staging check
  \item Full 2-stage simulation with realistic Falcon 9 parameters
  \item Plot: mass vs time (with visible discontinuous jumps at separation)
  \item Terminal output: Tsiolkovsky ideal $\Delta v$ vs simulated $\Delta v$ — compute both losses
  \item \textbf{Challenge:} Add a ``coasting'' phase between stages (engine-off glide before stage 2 ignition)
\end{enumerate}
\end{theorybox}

\end{document}
